\scanpage{75}
Jacobienne locale \inserted{\unsure{relative}} d'une localité \inserted{(régulière)} $A$ sur un corps $k$

$J_{A/k}$ est défini par le foncteur $\operatorname{Hom}_{k\text{-grp}}(J_{A/k},G)$ en $G$ (grp. \inserted{loc.} de type fini sur $k$) on doit avoir
\[
\operatorname{Hom}_{k\text{-}\unsure{gr|grp}}(J_{A/k},G)\simeq H_A^i(G)\qquad i=\dim A
\]
% slanted, upper left
\margin{Le foncteur en $G$ est \unsure{représentable}, \unsure{grâce au fait que} $H_A^j(G)=0$ si $j\unsure{<|\leq}i$ \unsure{lorsque|dès que} $A$ est \uuline{régulier}.}
$H_A^i$ est le foncteur cohomologique de localisation pour l'anneau local $A$ [$\text{de }\unsure{\operatorname{Spec}|\operatorname{Sp}}(A):\operatorname{Spec}(A)-{\unsure{m|\mathfrak m}}_A$], foncteur qui est défini à priori sur les faisceaux de Weil sur $\operatorname{Spec}(A)$. On a
% slanted, middle left
\margin{\unsure{Tous les} $J_{A/k}$ \unsure{sont des} proalgébriques \uline{\unsure{simples|stricts}} sur $k$.}
\[
\begin{aligned}
H_A^0(G)&=G(K)&&\text{si }A=K\text{ de dim }0\\
H_A^1(G)&=G(K)/G(A)&&\text{si }A\text{ de dim }1,\text{ de corps de fractions }K\\
H_A^i(G)&=H^{i-1}(\operatorname{Spec}(A)-m_A,G)&&\text{si }A\text{ de dim}\geq2
\end{aligned}
\]
\unsure{Je trouve|Je résume} :

\noindent a) $J_{K/k}$ est une \uline{extension de} $\mathbb Z_{(k)}$ (qui correspond au foncteur $G\longmapsto G(k)$) par un groupe connexe, qui en général a \uline{une partie abélienne} (c'est le \unsure{Albanese} de $K/k$), \uline{une partie multiplicative} \inserted{$M$}, en plus de la partie unipotente \inserted{$U$}.
% slanted, right; marginal leader, relative geometry only
\figure{\begin{tikzpicture}[x=1cm,y=1cm]
\fill (0.12,1.5) circle (0.025);
\draw (0.12,1.5) .. controls (-0.2,1.25) and (0.2,0.3) .. (0.55,-0.7) .. controls (0.68,-1.05) and (0.68,-1.5) .. (0.38,-1.52);
\end{tikzpicture}}
\margin{\unsure{Ces} \deleted{\unsure{J}} J sont \unsure{rel.t|relt.} projectifs pour les \unsure{ext.|cat.} \unsure{qui sont des} fibrés loc. triv.}
\inserted{$J_{K/k}$ \unsure{n'est pas projectif en général, mais son} $\operatorname{Ext}^1$ \unsure{avec un groupe} \uline{linéaire} \uline{connexe} est \unsure{tjrs|toujours} \unsure{trivial}.}
% slanted, lower left, first note
\margin{$\unsure{D^{(M)}|D(M)}$. N.B. $\operatorname{Hom}_{k\text{-grp}}(M,G_m)=\deleted{\unsure{K^*/k^*}}=$ grp des diviseurs \deleted{\unsure{de} \illegible{2}} \unsure{alg. équiv. à} $0$ (sur un modèle \unsure{régulier}).}
% slanted, lower left, second note
\margin{$\operatorname{Hom}(U,G_a)=K/k$ \deleted{\unsure{différent.} \illegible{3} \unsure{triv.}} \unsure{mais les} \unsure{mêmes} \unsure{sur|en} ${\unsure{\mathscr F|F}}_{\unsure{n-1|i-1}}(\unsure{k_0|k.})$ \unsure{et|de}.}

\noindent b) Si $A$ est de dim $1$, $J_{A/k}$ \uline{connexe}, \uline{linéaire} (i.e. plus de partie abélienne) mais \unsure{elle a} \deleted{\unsure{le plus souvent|la plus souvent}} une \deleted{\unsure{grosse}} partie multiplicative \inserted{$M\simeq G_m$}
% archive sheet number: 74
